\section*{Lectures en attente des scans}
\begingroup\small
Ces notes conservent les alternatives non résolues sans modifier la lecture retenue. Elles sont distinctes de l’apparat numéroté des différences entre textes diplomatique et corrigé. Chaque point nécessite le scan.

\medskip\noindent\textbf{SOL-EN-010 (folio 96).}\par
Lecture conservée : $f_{2*}(4)\circ(2)$. Proposition : $(f_{2*}(4)\circ(2))^{-1}$.
Les deux couches françaises fournies conservent le second chemin sans inverse dans le commentaire suivant 3.5. La traduction suit cette lecture. L’inverse proposé est une nouvelle correction mathématique ; la page imprimée et le diagramme complet doivent être revus avant de l’adopter.
\par\noindent\textit{Source: }\texttt{port/SOL\_NOTES\_5.md}. Statut : nécessite le scan.

\medskip\noindent\textbf{SOL-EN-017 (folio 144).}\par
Lecture conservée : $(\{y\}\times Y,B)$. Proposition : $(\{x\}\times Y,B)$.
La source française conserve $y$ dans ce couple, et l’anglais le reproduit. La proposition portant $x$ reste signalée pour contrôle de la page imprimée.
\par\noindent\textit{Source: }\texttt{port/SOL\_NOTES\_5.md}. Statut : nécessite le scan.

\medskip\noindent\textbf{SOL-EN-026 (folio 160).}\par
Lecture conservée : le localisé strict en $\bar z$. Proposition : le localisé strict en $\bar s$.
Les deux couches françaises conservent la localisation en $\bar z$ dans cette phrase. Le remplacement par $\bar s$ reste à vérifier sur la page imprimée.
\par\noindent\textit{Source: }\texttt{port/SOL\_NOTES\_5.md}. Statut : nécessite le scan.

\medskip\noindent\textbf{SOL-EN-028 (folio 161).}\par
Lecture conservée : $R\Gamma(B'\cup D,i'_!(E')|B'-(B'\cup D))$. Proposition : $R\Gamma(B'-D,i'_!(E')|B'-D)$.
Les deux couches françaises conservent l’union et la différence emboîtée de cette formule. La lecture enregistrée reste dans le texte ; la proposition portant $B'-D$ attend le contrôle de la page.
\par\noindent\textit{Source: }\texttt{port/SOL\_NOTES\_5.md}. Statut : nécessite le scan.

\medskip\noindent\textbf{SOL-EN-029 (folio 170).}\par
Lecture conservée : $\underline{\operatorname{Hom}}_A^\bullet(E,E)\to A$. But proposé : $A_\natural$.
La source française conserve $A$ dans cette phrase, tandis que les constructions voisines emploient $A_\natural$. La page imprimée est nécessaire avant d’adopter le changement proposé du but.
\par\noindent\textit{Source: }\texttt{port/SOL\_NOTES\_5.md}. Statut : nécessite le scan.

\medskip\noindent\textbf{SOL-EN-033 (folio 171).}\par
Renvoi conservé : 5.3.6. Renvoi proposé : 5.3.5.
Les deux couches françaises donnent 5.3.6 après (5.9.1). Le renvoi imprimé enregistré est conservé ; 5.3.5 reste une proposition à contrôler sur la page.
\par\noindent\textit{Source: }\texttt{port/SOL\_NOTES\_5.md}. Statut : nécessite le scan.

\medskip\noindent\textbf{SOL-EN-046 (folio 185).}\par
Lecture conservée : $E_2=KA_{X_1}$. Proposition : $F_1=KA_{X_1}$.
La couche française corrigée fournie porte elle aussi $E_2$. La proposition $F_1$ est enregistrée séparément, sans modifier cette lecture avant contrôle de l’imprimé.
\par\noindent\textit{Source: }\texttt{port/SOL\_NOTES\_5.md}. Statut : nécessite le scan.

\medskip\noindent\textbf{SOL-EN-050 (folio 191).}\par
Lecture conservée : $D_c({}_AX)=D_c({}_{A[G]}X)$. Premier membre proposé : $D_c(G,{}_AX)$.
Les deux couches françaises conservent le premier membre sans argument $G$. La lecture de la source reste dans le corps du texte ; l’ajout proposé est soumis à la vérification de la page.
\par\noindent\textit{Source: }\texttt{port/SOL\_NOTES\_5.md}. Statut : nécessite le scan.

\medskip\noindent\textbf{SOL-EN-052 (folio 195).}\par
Lecture conservée : $F^n\Lambda_X=0$ pour $n\geq1$. Borne proposée : $n\geq2$.
Les deux couches françaises conservent la borne $n\geq1$ à côté de la définition distincte de $F^1$. La lecture imprimée enregistrée est gardée et la borne proposée attend le scan.
\par\noindent\textit{Source: }\texttt{port/SOL\_NOTES\_5.md}. Statut : nécessite le scan.

\medskip\noindent\textbf{SOL-EN-053 (folio 195).}\par
But conservé : $c_2^!(\Lambda_{X-U,X})$. But proposé : $c_2^!(\Lambda_{U,X})$.
Les deux couches françaises donnent le premier de ces buts dans la correspondance graduée de degré un. Ce but est conservé ; le second reste une lecture proposée à vérifier.
\par\noindent\textit{Source: }\texttt{port/SOL\_NOTES\_5.md}. Statut : nécessite le scan.

\medskip\noindent\textbf{SOL-EN-056 (folio 197).}\par
Lecture conservée : tout point $x$ de $G$. Proposition : tout point $x$ de $C$.
Le français fourni porte lui aussi « pour tout point $x$ de $G$ ». Cette lecture reste dans les lecteurs ; $C$ est signalé comme alternative à vérifier sur l’imprimé.
\par\noindent\textit{Source: }\texttt{port/SOL\_NOTES\_5.md}. Statut : nécessite le scan.

\medskip\noindent\textbf{SOL-EN-057 (folio 197).}\par
Lecture conservée : $f_*M$. Proposition : $f_*L$.
Les deux couches françaises conservent $f_*M$ dans la proposition 6.19. La proposition portant $L$ est maintenue dans cette note pour contrôle de la page.
\par\noindent\textit{Source: }\texttt{port/SOL\_NOTES\_5.md}. Statut : nécessite le scan.

\medskip\noindent\textbf{SOL-EN-062 (folio 206).}\par
Lecture conservée : $\underline{\operatorname{Hom}}_A(M,\operatorname{Hom}_{C}(X,Y))$. Symbole extérieur proposé : $\operatorname{Hom}_A$.
Le français fourni souligne le Hom extérieur à cet endroit, comme le fait l’anglais porté. Cette notation explicite de la source reste inchangée ; le Hom sans soulignement attend le contrôle de l’imprimé.
\par\noindent\textit{Source: }\texttt{port/SOL\_NOTES\_5.md}. Statut : nécessite le scan.

\medskip\noindent\textbf{SOL-EN-073 (folio 371).}\par
Ponctuation conservée : « cf. SGA 3 I 4.7.2). » Proposition : « cf. SGA 3 I 4.7.2. »
Les deux couches françaises contiennent cette parenthèse fermante sans parenthèse ouvrante correspondante. La ponctuation enregistrée est conservée jusqu’au contrôle de la page imprimée.
\par\noindent\textit{Source: }\texttt{port/SOL\_NOTES\_5.md}. Statut : nécessite le scan.

\medskip\noindent\textbf{SOL-EN-075 (folio 375).}\par
Lecture conservée : formule de Weil 6.1 et invariants locaux au n\textsuperscript{o}~5. Proposition : formule 5.2 et invariants locaux au n\textsuperscript{o}~6.
Les deux couches françaises conservent ces deux renvois. La renumérotation proposée reste dans les notes en attente de contrôle, sans changer les renvois du corps du texte.
\par\noindent\textit{Source: }\texttt{port/SOL\_NOTES\_5.md}. Statut : nécessite le scan.

\medskip\noindent\textbf{SOL-EN-076 (folio 379).}\par
Lecture conservée : défini dans 4.1. Proposition : défini dans 3.4.
La proposition 3.6 française dit explicitement « défini dans 4.1 ». Le renvoi de la source est conservé ; 3.4 est enregistré comme alternative à vérifier.
\par\noindent\textit{Source: }\texttt{port/SOL\_NOTES\_5.md}. Statut : nécessite le scan.

\medskip\noindent\textbf{SOL-EN-077 (folio 380).}\par
Renvoi conservé : (4.3.1). Renvoi proposé : (3.7.1).
Les deux couches françaises donnent (4.3.1) dans la phrase suivant (3.7.2). Le renvoi conservé et l’alternative sont présentés ensemble pour la vérification de la page.
\par\noindent\textit{Source: }\texttt{port/SOL\_NOTES\_5.md}. Statut : nécessite le scan.

\medskip\noindent\textbf{SOL-EN-086 (folio 475).}\par
Lecture conservée : un endomorphisme de $K'$. Proposition : un endomorphisme de $K'_\nu$.
Le français fourni porte lui aussi $K'$ sans indice $\nu$ dans cette phrase. La lecture enregistrée reste dans le texte ; l’alternative indexée attend le scan.
\par\noindent\textit{Source: }\texttt{port/SOL\_NOTES\_5.md}. Statut : nécessite le scan.

\medskip\noindent\textbf{REVIEW-5-01 (folio 299).}\par
Lecture conservée : « faisceau normal ». Insertion manuscrite proposée : « faisceau conormal ».
Le dossier classe le mince ajout « co » comme une annotation du lecteur et conserve « normal ». Le même énoncé donne explicitement $N=J/J^2$ et emploie son dual. La formule est donc transcrite sans changer le nom enregistré. L’image de la page 299 n’est pas fournie : la nature de l’ajout manuscrit reste à vérifier.
\par\noindent\textit{Source: }\texttt{fixes/REVIEW\_5\_OPEN.jsonl}. Statut : nécessite le scan.

\medskip\noindent\textbf{GLOBAL-85-LABELS (folio Exposé V).}\par
Lecture conservée : les intitulés en gras et en casse mixte de la source. Alternative imprimée enregistrée : DEFINITION 1.1.~- et les intitulés en capitales correspondants.
Le dossier donne des exemples, sans fournir le relevé complet de chaque intitulé. Les scans correspondants manquent. Les formes de la source finale sont donc gardées ; aucun remplacement uniforme n’est déduit de ces seuls exemples.
\par\noindent\textit{Source: }\texttt{conventions/GLOBAL\_ISSUES\_SGA.md}. Statut : nécessite le scan.

\medskip\noindent\textbf{GLOBAL-85-ARROWS (folio 197-199).}\par
Lecture conservée : étiquettes des flèches horizontales inférieures placées au-dessous. Alternative imprimée enregistrée : étiquettes placées au-dessus.
Cette observation de placement n’a pas été détaillée diagramme par diagramme. Les sources, buts et étiquettes des flèches sont conservés. Les déplacer par une règle générale exige les pages absentes et une vérification de chaque diagramme.
\par\noindent\textit{Source: }\texttt{conventions/GLOBAL\_ISSUES\_SGA.md}. Statut : nécessite le scan.

\medskip\noindent\textbf{GLOBAL-89a (folio Exposé VII).}\par
Lecture conservée : intitulés en gras et en casse mixte. Alternatives imprimées enregistrées : capitales avec « .- », un tiret seul, ou ni point ni tiret selon l’énoncé.
Le dossier signale plusieurs ponctuations, dont un intitulé sans point ni tiret. Il ne donne pas de relevé complet permettant un remplacement uniforme. Les intitulés restent ceux de la source finale jusqu’à l’examen des pages.
\par\noindent\textit{Source: }\texttt{conventions/GLOBAL\_ISSUES\_SGA.md}. Statut : nécessite le scan.

\medskip\noindent\textbf{GLOBAL-89b (folio 281-295).}\par
Lecture conservée : les choix détaillés de la source à chaque occurrence. Alternatives restantes aux endroits non détaillés : $R$ ordinaire ou $\mathbb{R}$ à double trait ; le dossier nomme des exceptions dans $Rp_*$, $R_!p$, $Rf_*$ dans la preuve de 1.2.
Les contrôles par page ont déjà appliqué des corrections de caractère et ont expressément nommé des exceptions. Ces lectures et leur portée locale sont conservées. Les occurrences encore non examinées ne sont pas déduites d’une définition globale de macro.
\par\noindent\textit{Source: }\texttt{conventions/GLOBAL\_ISSUES\_SGA.md}. Statut : nécessite le scan.

\medskip\noindent\textbf{GLOBAL-89c (folio 288-289).}\par
Lecture conservée : $\mathbb{P}$ et $\mathbb{E}$. Alternatives imprimées enregistrées : $P$ et $E$ dactylographiés, en 21 endroits signalés.
La note globale n’individualise pas ces 21 occurrences. Les pages absentes sont nécessaires pour relier chaque constat au symbole précis. Les formes de la source finale sont conservées et les deux caractères restent signalés.
\par\noindent\textit{Source: }\texttt{conventions/GLOBAL\_ISSUES\_SGA.md}. Statut : nécessite le scan.

\medskip\noindent\textbf{GLOBAL-89e (folio Exposé VII).}\par
Lecture conservée : les capitales vérifiées à chaque occurrence. Alternatives restantes aux endroits non détaillés : Module et module.
Le dossier demande expressément un contrôle par occurrence, sans remplacement aveugle. Les formes finales déjà vérifiées sont conservées ; le constat général reste signalé sans modifier les passages non examinés.
\par\noindent\textit{Source: }\texttt{conventions/GLOBAL\_ISSUES\_SGA.md}. Statut : nécessite le scan.

\medskip\noindent\textbf{GLOBAL-89f (folio Exposé VII, notamment 295).}\par
Lecture conservée : les occurrences de $\mathcal{O}$ de la source. Alternatives imprimées enregistrées : $O$ dactylographié ordinaire et $\mathcal{O}_C$ manuscrit de forme cursive.
La note décrit plusieurs caractères sans donner de liste complète des occurrences. Un remplacement général effacerait cette variation enregistrée. Les caractères de la source finale restent en place ; les alternatives attendent les images.
\par\noindent\textit{Source: }\texttt{conventions/GLOBAL\_ISSUES\_SGA.md}. Statut : nécessite le scan.

\medskip\noindent\textbf{GLOBAL-VIII-CATEGORIES (folio 351-371).}\par
Lecture conservée : lettres cursives $\mathcal{C},\mathcal{D},\mathcal{A},\mathcal{B}$. Alternatives imprimées enregistrées : capitales dactylographiées $C,D,A,B$.
La note non numérotée pose une question de caractère, sans liste d’occurrences, et renvoie à la convention générale. Le texte final est conservé. Les quatre macros de catégories ne sont pas redéfinies globalement sur la seule base de cette note.
\par\noindent\textit{Source: }\texttt{conventions/GLOBAL\_ISSUES\_SGA.md}. Statut : nécessite le scan.

\medskip\noindent\textbf{GLOBAL-VIII-QUOTES (folio 351-371).}\par
Lecture conservée : les guillemets de chaque passage dans la source finale. Alternatives enregistrées : guillemets dactylographiés droits et guillemets français ; six paires sont signalées.
La note indique que les deux paires des pages 354--355 ont fait l’objet de corrections détaillées. Les formes finales sont conservées à ces endroits et ailleurs. Les images manquantes empêchent une nouvelle transcription uniforme des six paires.
\par\noindent\textit{Source: }\texttt{conventions/GLOBAL\_ISSUES\_SGA.md}. Statut : nécessite le scan.

\medskip\noindent\textbf{GLOBAL-VII-SET-DIFFERENCE (folio 301 et Exposé VII).}\par
Lecture conservée : les deux signes détaillés $X\mathbin{\dot{-}}Y$ et $X'\mathbin{\dot{-}}Y'$ à la page 301, ainsi que les signes existants ailleurs. Alternatives restantes ailleurs : moins pointé et signe ordinaire de différence ensembliste.
La source finale écrit explicitement le moins pointé dans les deux complémentaires de la page 301. Les autres occurrences n’ont pas été examinées dans la note fournie. Leurs signes restent donc inchangés jusqu’au contrôle des pages.
\par\noindent\textit{Source: }\texttt{conventions/GLOBAL\_ISSUES\_SGA.md}. Statut : nécessite le scan.

\medskip\noindent\textbf{GLOBAL-98b (folio 441).}\par
Lecture conservée dans les deux branches : $\underline{\operatorname{\mathbb{F}r}}_{Y/S}:Y\longrightarrow Y^{(p/S)}$. Alternative non résolue : $\underline{\operatorname{Fr}}_{Y/S}:Y\longrightarrow Y^{(p/S)}$.
Le point 98 laisse expressément cette occurrence ouverte face à la convention générale du point 78. L’entrée 572 du registre MISPRINTS\_5 décrit le F à double trait et son soulignement, tout en réservant son traitement à la décision sur les caractères du typographe. Le scan de la page 441 n’est pas fourni. La règle retenue pour les lectures non vérifiées impose donc de conserver la lecture imprimée enregistrée dans les six lecteurs ; l’alternative à F ordinaire reste ici en attente du scan. Le morphisme de Frobenius relatif $Y\to Y^{(p/S)}$ et les autres symboles de cette formule sont conservés.
\par\noindent\textit{Source: }\texttt{conventions/GLOBAL\_ISSUES\_SGA.md; fixes/MISPRINTS\_5.json, entry 572}. Statut : nécessite le scan.
\endgroup
\clearpage
